\chapter{Conclusões e Implementações Futuras}
\label{chap:conclusoes}

O modelo atual de otimização é projetado para melhorar a gestão do sistema de abastecimento de água. Contudo, a adição de novos elementos ao modelo pode enriquecê-lo e gerar resultados mais significativos. Uma proposta é a expansão da estratégia tarifária, que envolve a ampliação do modelo para incorporar uma abordagem mais detalhada das tarifas energéticas, que frequentemente variam em horários não-inteiros. Além disso, o modelo pode ser aprimorado para refletir diferentes bandeiras tarifárias. O resultado esperado dessa expansão é uma maior flexibilidade e precisão na gestão de custos, permitindo a adaptação das operações de bombeamento a um espectro mais amplo de condições tarifárias e alinhando a gestão de custos com as variações reais do mercado.

Outro aspecto crucial é a realização de uma análise de sensibilidade aprimorada. Para isso, a aplicação do Método de Monte Carlo pode ser explorada para avaliar o impacto de variações estocásticas nos custos de energia, padrões de consumo e outras variáveis operacionais. Este método permite simular diversos cenários, proporcionando uma visão mais clara das incertezas e riscos associados. O resultado esperado desta análise é a identificação de pontos críticos do sistema e uma capacidade aprimorada de prever e gerenciar riscos operacionais, oferecendo insights valiosos para o sistema em condições adversas.

Adicionalmente, a aplicação de outras técnicas de otimização é recomendada. Isso inclui a implementação de métodos como otimização estocástica, metaheurísticas ou novas políticas de gestão de consumo, para avaliar o desempenho do sistema. O resultado esperado dessas técnicas é um aumento na robustez do modelo, tornando-o mais resiliente a incertezas e variações no consumo e fornecimento de água.

Por fim, a expansão do conjunto de dados utilizado no modelo é fundamental. Enriquecer o modelo com dados mais abrangentes, incluindo séries temporais mais extensas, dados demográficos e informações econômicas, pode proporcionar uma compreensão mais profunda das tendências de longo prazo e das influências econômicas no consumo de água. O resultado esperado dessa expansão é a capacidade de realizar previsões e planejamentos mais precisos.